\documentclass[11pt]{amsart}
\usepackage[margin=1.1in]{geometry}
\usepackage{amsmath,amssymb,amsthm,mathtools}
\usepackage{booktabs}
\usepackage{array}
\usepackage[colorlinks,linkcolor=black,citecolor=black,urlcolor=blue]{hyperref}
\usepackage{url}

\newtheorem{theorem}{Theorem}[section]
\newtheorem{lemma}[theorem]{Lemma}
\newtheorem{proposition}[theorem]{Proposition}
\newtheorem{corollary}[theorem]{Corollary}
\theoremstyle{definition}
\newtheorem{definition}[theorem]{Definition}
\newtheorem{notation}[theorem]{Notation}
\theoremstyle{remark}
\newtheorem{remark}[theorem]{Remark}

\newcommand{\N}{\mathbb{N}}
\newcommand{\Q}{\mathbb{Q}}
\newcommand{\R}{\mathbb{R}}
\newcommand{\kk}{\Bbbk}
\newcommand{\g}{\mathfrak{g}}
\newcommand{\h}{\mathfrak{h}}
\newcommand{\lie}{\mathfrak{l}}
\newcommand{\kr}{\mathfrak{k}}
\newcommand{\ktop}{\mathfrak{k}^{\mathrm{top}}}
\newcommand{\af}{\mathfrak{a}}
\newcommand{\Lf}{\mathcal{L}}
\newcommand{\Ue}{\mathcal{U}}
\newcommand{\ad}{\operatorname{ad}}
\newcommand{\lead}{\operatorname{lead}}
\newcommand{\Hs}{\mathrm{H}}
\newcommand{\Sym}{\operatorname{Sym}}
\newcommand{\Irr}{\operatorname{Irr}}

\title[Relations in a Lie algebra of Tao]{The relations of a Lie subalgebra considered by Tao: a complete presentation}
\author{Qiyuan Ma}
\date{}
\makeatletter
\def\@setauthors{\begingroup\centering\normalfont
{\Large Qiyuan Ma\par}\vspace{2pt}
{\large Shanghai Jiao Tong University\par}\vspace{2pt}
{\large\href{mailto:mqyqwerty@sjtu.edu.cn}{mqyqwerty@sjtu.edu.cn}\par}\vspace{2pt}
{\large ORCID: \href{https://orcid.org/0009-0000-9532-0821}{0009-0000-9532-0821}\par}\vspace{10pt}
{\Large November 16, 2026\par}\vspace{2pt}
{\large Research preprint, version 1.1\par}
\endgroup}
\makeatother

\begin{document}

\begin{abstract}
Let $\g$ be the Lie algebra over a field of characteristic $0$ generated by $D_i,E_i$ ($i\ge 1$) subject to $[D_i,D_j]=[E_i,E_j]=0$ and $[D_i,E_j]=[D_j,E_i]$. Let $\h\subset\g$ be the subalgebra generated by the elements $[D_i,E_j]$ and $E_i$. On MathOverflow, T.~Tao asked whether the obvious relations $[D_i,E_j]=[D_j,E_i]$, $[E_i,E_j]=0$ and $[E_i,[D_j,E_k]]=[E_k,[D_j,E_i]]$ generate all relations of $\h$. We prove that they do. The proof has four parts. First, a polarisation argument reduces the problem to multilinear components. Second, we give a presentation of the ideal of $\g$ generated by the $E_i$. Third, a filtration argument shows that injectivity follows from an equality of dimensions; the associated graded algebra is computed explicitly. Fourth, the required dimensions are those of the multilinear components of $\g$, which are given by Eulerian numbers. These dimensions are equivalent to known results: the multilinear components of the enveloping algebra of the ideal generated by the $E_i$ coincide with the homology of the Losev--Manin spaces, whose Betti numbers were computed by Losev and Manin and whose algebra structure was presented by Dotsenko, Shadrin and Tamaroff. We explain this relation precisely, and we also include an independent, self-contained derivation, based on a Poincar\'e--Birkhoff--Witt basis of $U(\g)$ obtained from Bergman's diamond lemma and a sign-reversing involution. As by-products, $\h$ is a semidirect product of a free Lie algebra and an abelian Lie algebra, and the elements $[D_i,E_j]$ ($i\le j$) freely generate a free Lie subalgebra of $\g$. We were unable to find the affirmative answer to Tao's question or this freeness statement in the literature. We also report independent computer verifications.
\end{abstract}

\maketitle

\section{Introduction}

Throughout, $\kk$ is a field of characteristic $0$ (for instance $\R$), and $\N=\{1,2,3,\dots\}$. The following question was asked by T.~Tao on MathOverflow in February 2024 \cite{MO}.

Let $\g$ be the Lie algebra over $\kk$ generated by $D_i,E_i$, $i\in\N$, subject to
\begin{equation}\label{eq:g}
[D_i,D_j]=[E_i,E_j]=0,\qquad [D_i,E_j]=[D_j,E_i]\qquad(i,j\in\N).
\end{equation}
Let $\h\subset\g$ be the Lie subalgebra generated by the elements $[D_i,E_j]$ and $E_i$. These elements satisfy
\begin{gather}
[D_i,E_j]=[D_j,E_i],\qquad [E_i,E_j]=0, \tag{$*$}\label{eq:star}\\
[E_i,[D_j,E_k]]=[E_k,[D_j,E_i]], \tag{$**$}\label{eq:starstar}
\end{gather}
the second line following from the Jacobi identity. Tao asked whether \eqref{eq:star} and \eqref{eq:starstar} generate all relations of $\h$.

In a comment on the question \cite{YCor}, YCor reformulated it as follows. Let $\lie$ be the Lie algebra with generators $F_{ij}$, $E_i$ ($i,j\in\N$) and relators
\begin{equation}\label{eq:l}
F_{ij}=F_{ji},\qquad [E_i,E_j]=0,\qquad [E_i,F_{jk}]=[E_k,F_{ij}].
\end{equation}
There is a natural homomorphism $\theta\colon\lie\to\g$, $F_{ij}\mapsto[D_i,E_j]$, $E_i\mapsto E_i$, whose image is $\h$. The question asks whether $\theta$ is injective.

\begin{theorem}\label{thm:main}
The homomorphism $\theta\colon\lie\to\g$ is injective. Equivalently, the relations \eqref{eq:star} and \eqref{eq:starstar}, together with the axioms of a Lie algebra, generate all relations among the elements $[D_i,E_j]$ and $E_i$ of $\g$.
\end{theorem}

Along the way we obtain the following structural information.

\begin{corollary}\label{cor:intro}
\begin{enumerate}
\item The Lie algebra $\h\cong\lie$ is a semidirect product $\Lf\rtimes\af_E$. Here $\af_E$ is the abelian Lie algebra with basis $(E_i)_{i\in\N}$. The algebra $\Lf$ is a free Lie algebra with free generators $Y_W=\ad(E_{j_1})\cdots\ad(E_{j_r})[D_a,E_b]$, indexed by finite multisets $W=\{j_1,\dots,j_r,a,b\}$ of size at least $2$; the element $Y_W$ depends only on $W$. (Corollary~\ref{cor:structure}.)
\item In particular the elements $[D_i,E_j]$, $i\le j$, freely generate a free Lie subalgebra of $\g$. This answers affirmatively a further question raised by YCor \cite{YCor}. (Corollary~\ref{cor:free}.)
\item Fix $n\ge 2$. Consider the span of the brackets in $\g$ in which every index $1,\dots,n$ occurs exactly once and $D$'s occur exactly $a$ times. This subspace has dimension equal to the Eulerian number $A(n-1,a-1)$, the number of permutations of $n-1$ letters with $a-1$ descents. In total its dimension is $(n-1)!$ (summed over $a$). (Corollary~\ref{cor:eulerian}.) This statement is equivalent to known results on the homology of Losev--Manin spaces \cite{LM,DST}; see Section~\ref{sec:prior}.
\end{enumerate}
\end{corollary}

\subsection*{Status of the question and relation to prior work}
As of November 16, 2026 the MathOverflow question has no answer. The comments on it contain the reformulation above and partial observations. We searched for a later resolution (MathOverflow, arXiv, zbMATH Open, Crossref, OpenAlex, OEIS and general web search; Google Scholar and MathSciNet were not searched directly). We were unable to find an answer to the question, nor the freeness statement of Corollary~\ref{cor:intro}(2), in the literature.

Part of the machinery is, however, known in an equivalent form. Write the defining relators of the ideal of $\g$ generated by the $E_i$ (Section~\ref{sec:k}) in its enveloping algebra and restrict to multilinear components. One obtains exactly the relators of the twisted associative algebra $\mathrm{tHyperCom}$ of Dotsenko, Shadrin and Tamaroff, which they show to be the homology of the Losev--Manin spaces \cite[Prop.~2.4.2]{DST}. The Betti numbers of these spaces are Eulerian numbers \cite[Thm.~2.3]{LM}. Consequently, the Hilbert series of $U(\g)$ and of $\g$ (Corollary~\ref{cor:Useries}, Proposition~\ref{prop:gseries}) and Corollary~\ref{cor:intro}(3) follow from these known results. Section~\ref{sec:prior} gives the precise statement. Our proofs of these facts (Sections~\ref{sec:pbw}--\ref{sec:hilbert}) are independent of \cite{LM,DST} and serve to keep the paper self-contained. The remaining results are Lemmas~\ref{lem:A} and~\ref{lem:B}, Theorem~\ref{thm:main}, and Corollaries~\ref{cor:structure} and~\ref{cor:free}. We were unable to find them in the literature, but we make no stronger claim about them.

\subsection*{Outline of the proof}
\begin{itemize}
\item In Section~\ref{sec:polar} a polarisation argument reduces Theorem~\ref{thm:main} to the multilinear components, where each index occurs at most once.
\item In Section~\ref{sec:k} we show that $\g\cong\kr\rtimes\af_D$, where $\kr$ is the ideal generated by the $E_i$. We give an explicit presentation of $\kr$ in which $\lie$ appears as the algebra defined by the ``small'' generators and relations.
\item In Section~\ref{sec:filtration} an \emph{excess} filtration of $\kr$ is introduced. Its associated graded algebra is dominated by an explicit Lie algebra $\ktop$ (Lemma~\ref{lem:A}). If $\dim\kr=\dim\ktop$ in a multilinear degree, then $\theta$ is injective there.
\item In Section~\ref{sec:ktop} we determine $\ktop$ completely (Lemma~\ref{lem:B}) and compute its Hilbert series.
\item In Sections~\ref{sec:pbw}--\ref{sec:hilbert} we compute the Hilbert series of $\g$, hence of $\kr$, independently. We use a quadratic Gr\"obner basis of $U(\g)$, obtained from Bergman's diamond lemma \cite{Ber}, and count normal words by a sign-reversing involution. A generating-function identity shows that the two Hilbert series agree, which proves the theorem. Alternatively, the Hilbert series of $\g$ can be deduced from \cite{LM,DST}; see Section~\ref{sec:prior}.
\end{itemize}
The proof is self-contained apart from Bergman's diamond lemma \cite[Theorem~1.2]{Ber} and the Poincar\'e--Birkhoff--Witt (PBW) theorem (which is derived from the diamond lemma in \cite[\S3]{Ber}). Section~\ref{sec:prior} discusses the relation to the work of Losev--Manin and Dotsenko--Shadrin--Tamaroff. Section~\ref{sec:computations} reports independent computer checks: exact rational arithmetic up to multilinear degree $8$, and a modular computation, which is rigorous over $\Q$, in degree $9$.

\section{Conventions and standard facts}\label{sec:conv}

\subsection{Multisets and gradings}
A \emph{multiset} $T$ of indices is a finitely supported function $\N\to\mathbb{Z}_{\ge0}$, and $|T|$ denotes its size. We write $T+i$ for $T$ with one more copy of $i$, and $A+B$ for the sum of multisets. A multiset is a \emph{set} if all multiplicities are $\le1$.

All Lie algebras below are generated by elements carrying an \emph{index content} (a multiset), and all defining relations are homogeneous for it. Hence they are graded by the monoid of multisets. A component of content $S$ is called \emph{multilinear} if $S$ is a set. We also use two further gradings, the $D$-degree and the $E$-degree: $D_i$ has bidegree $(1,0)$ and $E_i$ has bidegree $(0,1)$.

Every permutation $\pi$ of $\N$ induces automorphisms of all algebras considered ($D_i\mapsto D_{\pi(i)}$, and so on), because all defining relations are invariant under renaming indices. Such an automorphism maps the component of content $S$ isomorphically onto that of content $\pi(S)$. Consequently, the dimension of a multilinear component depends only on $|S|$ and on the bidegree.

\subsection{Multilinear Hilbert series}
Let $M$ be a graded vector space of this kind with finite-dimensional components. Its \emph{multilinear Hilbert series} is the formal power series
\[
\Hs_M(x)=\sum_{n\ge0}\frac{x^n}{n!}\sum_{a,b\ge0}\dim M_{[n]}^{(a,b)}\,y^az^b\ \in\ \Q[y,z][[x]],
\]
where $[n]=\{1,\dots,n\}$ and $M_{[n]}^{(a,b)}$ is the component of content $[n]$ and bidegree $(a,b)$. We only use spaces in which the dimension of the content-$S$ component depends only on $|S|$.

\begin{lemma}[exponential formula]\label{lem:exp}
Let $M,M_1,\dots,M_k$ be graded spaces as above, with $M_\emptyset=0$.
\begin{enumerate}
\item $\Hs_{M_1\otimes\cdots\otimes M_k}=\Hs_{M_1}\cdots\Hs_{M_k}$, where content and bidegree are additive on tensor products.
\item $\Hs_{T(M)}=1/(1-\Hs_M)$ for the tensor algebra $T(M)=\bigoplus_{k\ge0}M^{\otimes k}$.
\item $\Hs_{\Sym(M)}=\exp(\Hs_M)$ for the symmetric algebra.
\end{enumerate}
\end{lemma}

\begin{proof}
(1) A tensor of contents $S_1,\dots,S_k$ has content $S_1+\dots+S_k$, which is a set $S$ if and only if the $S_i$ are disjoint sets with union $S$. Hence $(M_1\otimes\cdots\otimes M_k)_S=\bigoplus_{S=S_1\sqcup\cdots\sqcup S_k}\bigotimes_i(M_i)_{S_i}$. Counting ordered decompositions of $[n]$ with prescribed block sizes gives multinomial coefficients, which is exactly the product of exponential generating functions.
(2) follows from (1) by summing over $k$.
(3) The multilinear part of $\Sym^k(M)=(M^{\otimes k})_{\mathfrak S_k}$ of content $S$ is the space of coinvariants of $\bigoplus_{(S_1,\dots,S_k)}\bigotimes_i M_{S_i}$, the sum running over ordered decompositions into nonempty blocks. The group $\mathfrak S_k$ permutes these summands freely. The coinvariants are therefore isomorphic to the sum over unordered set partitions, and their series is $\Hs_M^k/k!$.
\end{proof}

\subsection{Free Lie algebras and enveloping algebras}
We use the following standard facts. For a set $X$, $\Lf(X)$ denotes the free Lie algebra on $X$.
\begin{enumerate}
\item[(F1)] Every map $X\to\Lf(X)$ extends uniquely to a derivation of $\Lf(X)$. More generally, let $\varphi\colon\Lf(X)\to\mathfrak m$ be a homomorphism. A \emph{$\varphi$-derivation} is a linear map $\delta\colon\Lf(X)\to\mathfrak m$ with $\delta[u,v]=[\delta u,\varphi v]+[\varphi u,\delta v]$. Such a $\delta$ is determined by its values on $X$.
\item[(F2)] Let $\mathfrak m$ be a Lie algebra and $\Sigma\subset\mathfrak m$ a subset. Suppose $\mathfrak m$ is generated by a set $X$. Then the ideal generated by $\Sigma$ is spanned by the elements $[x_1,[x_2,\dots,[x_k,\sigma]\cdots]]$ with $k\ge0$, $x_p\in X$ and $\sigma\in\Sigma$. Indeed, the span $Z$ of these elements is stable under $\ad x$ for $x\in X$. The set of $u$ with $[u,Z]\subseteq Z$ is a subalgebra, since $\ad[u,v]=[\ad u,\ad v]$, and it contains $X$, hence equals $\mathfrak m$.
\item[(F3)] If $\mathfrak m=\Lf(X)/(R)$, then $U(\mathfrak m)=\kk\langle X\rangle/(R)$, where $R$ is regarded as a set of noncommutative polynomials via $[u,v]=uv-vu$. In particular $U(\Lf(X))=\kk\langle X\rangle$. Both sides represent the functor $A\mapsto\{$maps $X\to A$ killing $R\}$ on associative algebras.
\item[(F4)] (PBW theorem.) For every Lie algebra $\mathfrak m$, the ordered monomials in a totally ordered basis of $\mathfrak m$ form a basis of $U(\mathfrak m)$; see \cite[\S3]{Ber}. If $\mathfrak m$ is graded as above and the basis is chosen homogeneous, it follows that $U(\mathfrak m)\cong\Sym(\mathfrak m)$ as graded vector spaces. Hence, by Lemma~\ref{lem:exp}, $\Hs_{U(\mathfrak m)}=\exp\Hs_{\mathfrak m}$ whenever $\mathfrak m_\emptyset=0$. Applied to $\Lf(X)$ with graded $X$, (F3) and Lemma~\ref{lem:exp}(2) give $\Hs_{\Lf(X)}=-\log(1-\Hs_{\kk X})$.
\end{enumerate}
All identities between power series are identities in $\Q[y,z][[x]]$. Logarithms are taken of series with constant term $1$, and $\log(PQ)=\log P+\log Q$ holds for such series.

\section{Reduction to multilinear components}\label{sec:polar}

For a vector space $V$ we define Lie algebras functorially in $V$. Let $U=\kk d\oplus\kk e$, and put $D(v)=v\otimes d$ and $E(v)=v\otimes e$. Define
\begin{align*}
\g(V)&=\Lf(V\otimes U)\big/\big([D(v),D(w)],\ [E(v),E(w)],\ [D(v),E(w)]-[D(w),E(v)] : v,w\in V\big),\\
\lie(V)&=\Lf(V\oplus\Sym^2V)\big/\big([E(v),E(w)],\ [E(u),F(vw)]-[E(w),F(uv)] : u,v,w\in V\big).
\end{align*}
Here, in the second line, $E(v)$ is $v\in V$ and $F(vw)$ is the product $vw\in\Sym^2V$. Let $\theta_V\colon\lie(V)\to\g(V)$ be given by $E(v)\mapsto E(v)$ and $F(vw)\mapsto[D(v),E(w)]$. This is well defined on $\Sym^2V$ because $[D(v),E(w)]=[D(w),E(v)]$ in $\g(V)$. It respects the relations: the first is clear, and the second is mapped to
\[
[E(u),[D(v),E(w)]]-[E(w),[D(v),E(u)]]
=[[E(u),D(v)],E(w)]-[E(w),[D(v),E(u)]]=0,
\]
using Jacobi, $[E(u),E(w)]=0$, and $[E(u),D(v)]=-[D(v),E(u)]$.

Each linear map $\phi\colon V\to V'$ induces homomorphisms $\lie(\phi)$ and $\g(\phi)$, since the relations are defined for all vectors, and $\theta_{V'}\circ\lie(\phi)=\g(\phi)\circ\theta_V$.

For $V=\kk^{(\N)}$ with basis $(v_i)$, put $D_i=D(v_i)$, $E_i=E(v_i)$ and $F_{ij}=F(v_iv_j)$. Then $\g(V)=\g$ and $\lie(V)=\lie$, since the relations for all vectors follow by multilinearity from those for basis vectors. These are exactly \eqref{eq:g} and \eqref{eq:l}, and $\theta_V=\theta$. The content grading is the grading by the weights of the diagonal torus: if $\tau=\operatorname{diag}(t_i)$, then $\lie(\tau)$ and $\g(\tau)$ act on the component of content $\mu$ by $t^\mu=\prod t_i^{\mu(i)}$. Since $\theta$ is graded, $\ker\theta$ is a graded subspace.

\begin{proposition}[polarisation]\label{prop:polar}
If $\ker\theta$ has a nonzero component of content $\mu$ with $|\mu|=n$, then $\ker\theta$ has a nonzero component of content $[n]$.
\end{proposition}

\begin{proof}
Let $0\ne x\in(\ker\theta)_\mu$. Let $i_1,\dots,i_r$ be the distinct elements of the support of $\mu$, with multiplicities $m_1,\dots,m_r$, so $\sum m_j=n$. Partition $[n]=P_1\sqcup\cdots\sqcup P_r$ with $|P_j|=m_j$.

For $t=(t_1,\dots,t_n)\in\kk^n$ define linear maps $\beta_t,\alpha\colon V\to V$ by
\[
\beta_t(v_{i_j})=\sum_{p\in P_j}t_pv_p,\quad \beta_t(v_i)=0\ (i\notin\{i_1,\dots,i_r\});\qquad
\alpha(v_p)=v_{i_j}\ (p\in P_j),\quad \alpha(v_q)=0\ (q>n).
\]
Then $\alpha\beta_t(v_{i_j})=c_jv_{i_j}$ with $c_j=\sum_{p\in P_j}t_p$, and $\alpha\beta_t(v_i)=0$ for the other $i$.

The element $x$ is a linear combination of brackets of generators whose indices lie in $\{i_1,\dots,i_r\}$, with total content $\mu$. Hence $\lie(\alpha\beta_t)(x)=\prod_jc_j^{m_j}\,x$.

On the other hand, $\beta_t=\tau_t\beta_1$, where $\tau_t$ is diagonal with entries $t_p$ on $v_p$ for $p\le n$ and $1$ elsewhere. Write $\lie(\beta_1)(x)=\sum_\nu x_\nu$ as a finite sum of homogeneous components; each $\nu$ is supported in $[n]$ and has size $n$. Then $\lie(\beta_t)(x)=\sum_\nu t^\nu x_\nu$, so
\[
\sum_\nu t^\nu\,\lie(\alpha)(x_\nu)=\prod_{j=1}^r\Big(\sum_{p\in P_j}t_p\Big)^{m_j}x\qquad\text{for all }t\in\kk^n.
\]
Both sides are polynomial maps $\kk^n\to\lie_\mu$, and $\kk$ is infinite, so the coefficients of $t_1t_2\cdots t_n$ agree. This gives $\lie(\alpha)(x_{[n]})=m_1!\cdots m_r!\,x\neq0$, using that the characteristic is $0$. In particular $x_{[n]}\neq0$.

Finally, $\theta(\lie(\beta_1)x)=\g(\beta_1)\theta(x)=0$ and $\theta$ is graded, so $x_{[n]}\in(\ker\theta)_{[n]}$.
\end{proof}

Hence Theorem~\ref{thm:main} follows once we show $(\ker\theta)_S=0$ for every finite set $S\subset\N$. By the symmetry under permutations of $\N$, it suffices to treat $S=[n]$.

\section{A presentation of the ideal generated by the \texorpdfstring{$E_i$}{E\_i}}\label{sec:k}

Let $\af_D$ be the abelian Lie algebra with basis $(D_i)_{i\in\N}$. Let $F=\Lf(G_T: T\text{ a nonempty multiset})$ be the free Lie algebra on symbols $G_T$; $G_T$ has content $T$ and bidegree $(|T|-1,1)$. By (F1), for each $k\in\N$ there is a unique derivation $\partial_k$ of $F$ with $\partial_kG_T=G_{T+k}$. The derivations $\partial_k$ pairwise commute, since $[\partial_k,\partial_l]$ is a derivation vanishing on the generators. For a multiset $S$ put $\partial_S=\prod_{s\in S}\partial_s$, and define
\[
r_{i,j,S}=\partial_S[G_i,G_j]\in F\qquad(i,j\in\N,\ S\text{ a multiset}).
\]
Explicitly, list $S=\{s_1,\dots,s_m\}$ with repetitions and write $S_P=\{s_p:p\in P\}$ for $P\subseteq[m]$. Then
\begin{equation}\label{eq:r}
r_{i,j,S}=\sum_{P\subseteq[m]}[G_{S_P+i},\,G_{S_{[m]\setminus P}+j}].
\end{equation}
Let $I\subset F$ be the ideal generated by all $r_{i,j,S}$, and put $\kr=F/I$. Since $\partial_kr_{i,j,S}=r_{i,j,S+k}$, and $\partial_k[u,\rho]=[\partial_ku,\rho]+[u,\partial_k\rho]$, each $\partial_k$ preserves $I$ and descends to a derivation of $\kr$. The descended derivations still commute. Therefore $D_k\mapsto\partial_k$ defines an action of $\af_D$ on $\kr$ by derivations, and we form the semidirect product $\kr\rtimes\af_D$, in which $[D_k,u]=\partial_ku$.

\begin{proposition}\label{prop:k}
\begin{enumerate}
\item In $\g$, the element $\ad(D_{t_1})\cdots\ad(D_{t_{m}})E_t$ depends only on the multiset $T=\{t_1,\dots,t_m,t\}$. Denote it by $\mathbf G_T$.
\item There is an isomorphism $\psi\colon\kr\rtimes\af_D\to\g$ with $\psi(G_T)=\mathbf G_T$ and $\psi(D_i)=D_i$. Its inverse $\varphi$ satisfies $\varphi(D_i)=D_i$ and $\varphi(E_i)=G_i$.
\item Under $\psi$, the ideal of $\g$ generated by the $E_i$ corresponds to $\kr$, and $\g=\psi(\kr)\oplus\af_D$ as graded vector spaces.
\item Let $F^{\mathrm{sm}}=\Lf(G_T:|T|\le2)\subset F$, and let $I_{\lie}\subset F^{\mathrm{sm}}$ be the ideal generated by the $r_{i,j,S}$ with $|S|\le1$. Then $\lie\cong F^{\mathrm{sm}}/I_{\lie}$ via $E_i\mapsto G_i$ and $F_{ij}\mapsto G_{\{i,j\}}$. Under this identification, $\theta$ equals $\psi\circ\iota$, where $\iota\colon F^{\mathrm{sm}}/I_{\lie}\to F/I=\kr$ is induced by the inclusion. In particular $\ker\theta=\ker\iota=(F^{\mathrm{sm}}\cap I)/I_{\lie}$.
\end{enumerate}
\end{proposition}

\begin{proof}
(1) Since $[D_i,D_j]=0$, the operators $\ad D_i$ commute, so the order of the $D$'s is irrelevant. It remains to exchange the distinguished index $t$ with some $t'$. Write $T=R+t+t'$. Then
\[
\ad(D_R)\ad(D_{t'})E_t=\ad(D_R)[D_{t'},E_t]=\ad(D_R)[D_t,E_{t'}]=\ad(D_R)\ad(D_t)E_{t'}.
\]

(2) Let $\psi_0\colon F\to\g$ be the homomorphism with $G_T\mapsto\mathbf G_T$. The maps $\psi_0\circ\partial_k$ and $\ad(D_k)\circ\psi_0$ are both $\psi_0$-derivations. They agree on generators, because $\psi_0(G_{T+k})=\mathbf G_{T+k}=[D_k,\mathbf G_T]$. By (F1) they are equal. Hence $\psi_0(r_{i,j,S})=\ad(D_S)[E_i,E_j]=0$, so $\psi_0$ factors through $\kr$, and the induced map intertwines $\partial_k$ with $\ad D_k$. Together with $D_i\mapsto D_i$ it therefore defines a homomorphism $\psi\colon\kr\rtimes\af_D\to\g$.

Conversely, define $\varphi$ on generators by $D_i\mapsto D_i$ and $E_i\mapsto G_i$. The relations \eqref{eq:g} are respected:
\begin{itemize}
\item $[D_i,D_j]\mapsto0$, since $\af_D$ is abelian;
\item $[E_i,E_j]\mapsto[G_i,G_j]=r_{i,j,\emptyset}=0$;
\item $[D_i,E_j]\mapsto\partial_iG_j=G_{\{i,j\}}$, which is symmetric in $i,j$.
\end{itemize}
Now $\psi\varphi$ is the identity on the generators $D_i,E_i$ of $\g$. Also $\varphi\psi(G_T)=\varphi(\ad(D_{T-t})E_t)=\partial_{T-t}G_t=G_T$, and $\varphi\psi(D_i)=D_i$. Hence $\psi$ and $\varphi$ are mutually inverse.

(3) The subspace $\kr$ is an ideal of $\kr\rtimes\af_D$, it contains $G_i=\varphi(E_i)$, and it is generated as an ideal by these elements, since $G_T=\partial_{T-t}G_t$. The quotient by it is $\af_D$.

(4) Identify $E_i=G_i$ and $F_{ij}=G_{\{i,j\}}$, so the generators $F_{ij}$ are indexed by $2$-element multisets and $F_{ij}=F_{ji}$ is built in. Then $r_{i,j,\emptyset}=[E_i,E_j]$, and by \eqref{eq:r}
\[
r_{i,j,\{s\}}=[G_{\{s,i\}},G_j]+[G_i,G_{\{s,j\}}]=[E_i,F_{sj}]-[E_j,F_{is}].
\]
This is the relator $[E_i,F_{jk}]-[E_k,F_{ij}]$ of \eqref{eq:l} with $(i,j,k)$ replaced by $(i,s,j)$. Thus the defining relators of $\lie$ and the generators of $I_{\lie}$ span the same space, and $\lie\cong F^{\mathrm{sm}}/I_{\lie}$. Finally, $\psi\iota(E_i)=E_i$ and $\psi\iota(F_{ij})=\mathbf G_{\{i,j\}}=[D_i,E_j]$, which agrees with $\theta$. Since $\psi$ is injective, $\ker\theta=\ker\iota$.
\end{proof}

In particular $\h=\theta(\lie)$ is the subalgebra of $\kr\cong\psi(\kr)$ generated by the $G_T$ with $|T|\le2$.

\section{The excess filtration}\label{sec:filtration}

Call $G_T$ \emph{small} if $|T|\le2$ and \emph{big} if $|T|\ge3$. Give $F$ the additional grading by \emph{excess}, defined on generators by
\[
\varepsilon(G_T)=\max(|T|-2,0)
\]
and extended additively to Lie monomials; this is a grading since $F$ is free. The weight-$0$ part $F_0$ is spanned by Lie monomials in small letters, so $F_0=F^{\mathrm{sm}}$ (the subalgebra of Proposition~\ref{prop:k}(4)). All gradings (content, bidegree and excess) are compatible, and every content component $F_\mu$ is finite-dimensional. For $0\ne u\in F$ let $\lead(u)$ denote its nonzero excess-homogeneous component of highest weight.

\begin{lemma}\label{lem:lead}
For a multiset $S$ with $|S|=m$, define
\[
\lambda_{i,j,S}=
\begin{cases}
[G_i,G_j] & m=0,\\
[G_{S+i},G_j]+[G_i,G_{S+j}] & m\ge1.
\end{cases}
\]
Then $\lambda_{i,j,S}$ is excess-homogeneous of weight $\max(m-1,0)$. If $i=j$, then $r_{i,j,S}=\lambda_{i,j,S}=0$. If $i\neq j$, then $\lambda_{i,j,S}\neq0$ and $\lead(r_{i,j,S})=\lambda_{i,j,S}$. If $m\le1$, then $r_{i,j,S}=\lambda_{i,j,S}$.
\end{lemma}

\begin{proof}
If $i=j$, then $r_{i,i,S}=\partial_S[G_i,G_i]=0$ and $\lambda_{i,i,S}=0$. Let $i\ne j$. The term of \eqref{eq:r} indexed by $P$ with $|P|=a$ has excess $\max(a-1,0)+\max(m-a-1,0)$. For $m\le1$ every term has excess $0$, and \eqref{eq:r} gives $r_{i,j,S}=\lambda_{i,j,S}$.

For $m\ge2$ the excess equals $m-1$ for $a\in\{0,m\}$ and $m-2$ for $1\le a\le m-1$. The two terms of top weight are those with $P=\emptyset$ and $P=[m]$, each occurring exactly once, and their sum is $\lambda_{i,j,S}$.

It remains to see that $\lambda_{i,j,S}\neq0$. For $m=0$, $[G_i,G_j]\ne0$ in the free Lie algebra since $i\ne j$. For $m\ge1$, the brackets $[G_{S+i},G_j]$ and $[G_i,G_{S+j}]$ are brackets of two distinct letters, and their sum vanishes only if $\{G_{S+i},G_j\}=\{G_i,G_{S+j}\}$. This is impossible: $G_{S+i}\neq G_i$ since $m\ge1$, and $G_{S+i}=G_{S+j}$ would force $i=j$.
\end{proof}

Let $J\subset F$ be the ideal generated by all $\lambda_{i,j,S}$, and put
\[
\ktop=F/J.
\]
Let $F_{\le p}=\bigoplus_{q\le p}F_q$, and define the associated graded of $I$ by $\operatorname{gr}I=\bigoplus_p\operatorname{gr}_pI$, where $\operatorname{gr}_pI=\pi_p(I\cap F_{\le p})\subseteq F_p$ and $\pi_p$ is the projection onto $F_p$.

\begin{lemma}\label{lem:grI}
\begin{enumerate}
\item $\operatorname{gr}I$ is an ideal of $F$, and $J\subseteq\operatorname{gr}I$.
\item For every content $\mu$, $\dim(\operatorname{gr}I)_\mu=\dim I_\mu$. Consequently $\dim\kr_\mu\le\dim\ktop_\mu$, with equality if and only if $J_\mu=(\operatorname{gr}I)_\mu$.
\end{enumerate}
\end{lemma}

\begin{proof}
(1) Let $u\in I\cap F_{\le p}$, and let $y$ be a Lie monomial of excess $q$. Then $[y,u]\in I\cap F_{\le p+q}$ and $\pi_{p+q}[y,u]=[y,\pi_pu]$. Hence $[y,\operatorname{gr}_pI]\subseteq\operatorname{gr}_{p+q}I$, so $\operatorname{gr}I$ is an ideal.

For each $(i,j,S)$ with $i\neq j$, put $p=\max(|S|-1,0)$. Lemma~\ref{lem:lead} gives $r_{i,j,S}\in I\cap F_{\le p}$ with $\pi_p(r_{i,j,S})=\lambda_{i,j,S}$. So every generator of $J$ lies in $\operatorname{gr}I$, and therefore $J\subseteq\operatorname{gr}I$.

(2) On the finite-dimensional space $I_\mu$, the filtration $I_\mu\cap F_{\le p}$ is finite. Its successive quotients are isomorphic to $(\operatorname{gr}_pI)_\mu$. Hence $\dim(\operatorname{gr}I)_\mu=\dim I_\mu$, and $\dim\kr_\mu=\dim F_\mu-\dim(\operatorname{gr}I)_\mu\le\dim F_\mu-\dim J_\mu=\dim\ktop_\mu$.
\end{proof}

\begin{lemma}\label{lem:J0}
The weight-$0$ part of $J$ is $J_0=J\cap F_0=I_{\lie}$. Consequently $\lie\cong F_0/J_0$ is the excess-$0$ part of $\ktop$.
\end{lemma}

\begin{proof}
By (F2), applied with $X$ the set of all $G_T$, the ideal $J$ is spanned by the excess-homogeneous elements $[y_1,[y_2,\dots,[y_k,\lambda]\cdots]]$, with $y_p$ letters and $\lambda=\lambda_{i,j,S}$. Such an element has excess $\sum\varepsilon(y_p)+\max(|S|-1,0)$. This is $0$ if and only if all $y_p$ are small and $|S|\le1$, and in that case $\lambda_{i,j,S}=r_{i,j,S}$. Since $J$ is excess-graded, $J_0$ is the span of these weight-$0$ elements. By (F2) applied inside $F_0$, this span is exactly $I_{\lie}$. The last statement follows from Proposition~\ref{prop:k}(4).
\end{proof}

\begin{lemma}[Lemma A]\label{lem:A}
Let $\mu$ be a content. If $\dim\kr_\mu=\dim\ktop_\mu$, then $(\ker\theta)_\mu=0$.
\end{lemma}

\begin{proof}
By Lemma~\ref{lem:grI}(2), $J_\mu=(\operatorname{gr}I)_\mu$. Let $u\in(F_0\cap I)_\mu$. Then $u\in I\cap F_{\le0}$, so $u=\pi_0(u)\in(\operatorname{gr}_0I)_\mu=(J_0)_\mu=(I_{\lie})_\mu$, by Lemma~\ref{lem:J0}. Hence $(F_0\cap I)_\mu=(I_{\lie})_\mu$, and $(\ker\theta)_\mu=0$ by Proposition~\ref{prop:k}(4).
\end{proof}

\section{The structure of \texorpdfstring{$\ktop$}{k-top}}\label{sec:ktop}

In $\ktop$ write $E_i=G_i$. By Lemma~\ref{lem:lead}, the defining relations of $\ktop$ are $[E_i,E_j]=0$ and, for $|S|\ge1$, $[G_{S+i},G_j]+[G_i,G_{S+j}]=0$. Setting $T=S+i$ (so $|T|\ge2$ and $t=i\in T$), the latter reads
\begin{equation}\label{eq:exchange}
[E_j,G_T]=[E_t,G_{T-t+j}]\qquad(|T|\ge2,\ t\in T,\ j\in\N).
\end{equation}
In words, $\ad(E_j)G_T$ is unchanged when $j$ is exchanged with any element of $T$. Since $[E_i,E_j]=0$, the operators $\ad E_i$ pairwise commute in $\ktop$.

Let $Y$ be the set of symbols $Y_{W,s}$, where $W$ is a multiset and $2\le s\le|W|$. The symbol $Y_{W,s}$ has content $W$, bidegree $(s-1,|W|-s+1)$ and excess $s-2$. Let $\af_E$ be abelian with basis $(E_i)$; it has excess $0$. By (F1) the rule $E_j\cdot Y_{W,s}=Y_{W+j,s}$ defines commuting derivations of $\Lf(Y)$. Form $\Lf(Y)\rtimes\af_E$.

\begin{lemma}[Lemma B]\label{lem:B}
\begin{enumerate}
\item For $W=C+T$ with $|T|=s\ge2$, write $C=\{j_1,\dots,j_r\}$. Then the element $\ad(E_{j_1})\cdots\ad(E_{j_r})G_T\in\ktop$ depends only on $(W,s)$.
\item There is an isomorphism $\beta\colon\Lf(Y)\rtimes\af_E\to\ktop$ with $\beta(E_i)=E_i$ and $\beta(Y_{W,s})=\ad(E_{C})G_T$ for any splitting as in (1). Its inverse $\alpha$ is given by $\alpha(E_i)=E_i$ and $\alpha(G_T)=Y_{T,|T|}$ for $|T|\ge2$. Both maps preserve content, bidegree and excess.
\end{enumerate}
\end{lemma}

\begin{proof}
(1) The order of the $\ad E_j$ is irrelevant, since they commute. Suppose $(C,T)$ and $(C',T')$ are two splittings of $W$ with $|T|=|T'|=s$ and $T\neq T'$. Choose $t\in T-T'$ and $j\in T'-T$ (multiset differences; both are nonempty since $|T|=|T'|$). The element $j$ has larger multiplicity in $T'\subseteq W$ than in $T$, hence occurs in $C=W-T$. Exchanging $j$ and $t$ gives a new splitting $(C-j+t,\,T-t+j)$, and $|T-T'|$ decreases by one. Each such exchange preserves the element: using commutativity, we may place $\ad E_j$ innermost, and then apply \eqref{eq:exchange}, valid since $|T|\ge2$:
\[
\ad(E_{C-j})[E_j,G_T]=\ad(E_{C-j})[E_t,G_{T-t+j}]=\ad(E_{C-j+t})G_{T-t+j}.
\]

(2) First, $\alpha$ is well defined. Define it on the generators of $F$ as stated. Then $[E_i,E_j]\mapsto0$. For \eqref{eq:exchange}, both sides map to $E_j\cdot Y_{T,|T|}=Y_{T+j,|T|}=E_t\cdot Y_{T-t+j,|T|}$. Hence $\alpha$ kills $J$.

Next, $\beta$ is well defined. Let $\beta_0\colon\Lf(Y)\to\ktop$ be defined freely by (1). Then $\beta_0\circ(E_j\cdot)$ and $\ad(E_j)\circ\beta_0$ are $\beta_0$-derivations, which agree on generators because
\[
\beta_0(Y_{W+j,s})=\ad(E_{C+j})G_T=[E_j,\beta_0(Y_{W,s})].
\]
By (F1) they are equal. Since the $E_i$ commute in $\ktop$, $\beta_0$ extends to a homomorphism $\beta$ of the semidirect product.

Both composites are the identity on generators:
\begin{itemize}
\item $\alpha\beta(Y_{W,s})=\alpha(\ad(E_{C})G_T)=E_{C}\cdot Y_{T,s}=Y_{W,s}$;
\item $\beta\alpha(G_T)=\beta(Y_{T,|T|})=G_T$;
\item $\alpha$ and $\beta$ fix each $E_i$.
\end{itemize}
The gradings are visibly preserved. For instance, $\ad(E_{C})G_T$ has excess $|T|-2=s-2$.
\end{proof}

\begin{corollary}\label{cor:ktopseries}
$\Hs_{\ktop}(x)=xz-\log\big(1-\mathcal Y(x)\big)$, where
\[
\mathcal Y(x)=\sum_{w\ge2}\frac{x^w}{w!}\sum_{s=2}^{w}y^{s-1}z^{w-s+1}.
\]
Moreover the excess-$0$ part satisfies $\Hs_{\lie}(x)=xz-\log\big(1-\mathcal Y_2(x)\big)$ with $\mathcal Y_2(x)=\sum_{w\ge2}\frac{x^w}{w!}yz^{w-1}$.
\end{corollary}

\begin{proof}
As a graded vector space, $\ktop\cong\Lf(Y)\oplus\af_E$. Multilinear generators $Y_{W,s}$ have $W$ a set. For $|W|=w$ they contribute $\sum_{s=2}^wy^{s-1}z^{w-s+1}$, so $\Hs_{\kk Y}=\mathcal Y$. Also $\Hs_{\af_E}=xz$. Now apply (F4).

The excess-$0$ part of $\Lf(Y)\rtimes\af_E$ is $\Lf(Y_{\cdot,2})\rtimes\af_E$, because the generators of excess $0$ are exactly the $Y_{W,2}$ and a Lie monomial has excess $0$ if and only if all its letters do. The subalgebra of a free Lie algebra generated by a subset of the free generators is free on that subset; it is a retract, by killing the other generators. Combining this with Lemma~\ref{lem:J0} gives the second formula.
\end{proof}

\section{A PBW basis of \texorpdfstring{$U(\g)$}{U(g)}}\label{sec:pbw}

By (F3), $U(\g)=\kk\langle X\rangle/(\mathcal R)$ with $X=\{D_i,E_i:i\in\N\}$. The ideal $(\mathcal R)$ is generated by the following elements, for $i<j$:
\[
D_iD_j-D_jD_i,\qquad E_iE_j-E_jE_i,\qquad D_iE_j-E_jD_i-D_jE_i+E_iD_j.
\]
Order $X$ by index first and then $D<E$:
\[
D_1<E_1<D_2<E_2<\cdots.
\]
Order words by length, and words of equal length lexicographically (deglex). This is a total order on the free monoid $\langle X\rangle$ satisfying $B<B'\Rightarrow ABC<AB'C$. It is a well-order, since for each length the lexicographic order on words over a well-ordered alphabet is a well-order. In particular it satisfies the descending chain condition.

We use the reduction system $\mathcal S$ consisting of the rules, for all $i<j$:
\begin{align*}
(\sigma_1)\quad & D_jD_i\ \to\ D_iD_j,\\
(\sigma_2)\quad & E_jE_i\ \to\ E_iE_j,\\
(\sigma_3)\quad & E_jD_i\ \to\ D_iE_j-D_jE_i+E_iD_j.
\end{align*}
The differences of the two sides are exactly the generators of $(\mathcal R)$, up to sign. In each rule, every word on the right is smaller than the word on the left. This holds because the first letter on the left is the largest of all first letters involved: $D_j>D_i$, $E_j>E_i$, and $E_j>D_i,D_j,E_i$. So the order is compatible with $\mathcal S$ in the sense of \cite[\S1]{Ber}.

Call a word of length $2$ \emph{bad} if it is the left side of a rule. Thus $X_jY_i$ is bad if and only if $j>i$ and $(X,Y)\neq(D,E)$. A word is \emph{normal} (irreducible) if it contains no bad subword.

\begin{proposition}\label{prop:ambig}
All ambiguities of $\mathcal S$ are resolvable.
\end{proposition}

\begin{proof}
All left sides have length $2$, and distinct rules have distinct left sides, so there are no inclusion ambiguities. An overlap ambiguity is a word $X_kY_jZ_i$ such that $X_kY_j$ and $Y_jZ_i$ are both bad. This forces $k>j>i$, $(X,Y)\neq(D,E)$ and $(Y,Z)\neq(D,E)$. Hence $(X,Y,Z)\in\{DDD,EDD,EED,EEE\}$.

The reductions below only compare indices, so it suffices to treat $(k,j,i)=(3,2,1)$; the general case follows by renaming $3,2,1$ to $k,j,i$. In each case we reduce the left pair, respectively the right pair, first, and then reduce to normal form. Words that are already normal are not reduced further; they are recognised by the criterion above.

\emph{$DDD$ and $EEE$.} Only $\sigma_1$ (resp.\ $\sigma_2$) applies. Both routes reach the sorted word:
\begin{align*}
D_3D_2D_1&\to D_2D_3D_1\to D_2D_1D_3\to D_1D_2D_3,\\
D_3D_2D_1&\to D_3D_1D_2\to D_1D_3D_2\to D_1D_2D_3,
\end{align*}
and likewise with $E$ in place of $D$.

\emph{$EDD$.} Left route, starting with $\sigma_3$ on $E_3D_2$:
\begin{align*}
E_3D_2D_1&\to D_2E_3D_1-D_3E_2D_1+E_2D_3D_1.
\end{align*}
We reduce each of the three terms:
\begin{align*}
D_2E_3D_1&\to D_2D_1E_3-D_2D_3E_1+D_2E_1D_3\to D_1D_2E_3-D_2D_3E_1+D_2E_1D_3,\\
-D_3E_2D_1&\to -D_3D_1E_2+D_3D_2E_1-D_3E_1D_2\to -D_1D_3E_2+D_2D_3E_1-D_3E_1D_2,\\
E_2D_3D_1&\to E_2D_1D_3\to D_1E_2D_3-D_2E_1D_3+E_1D_2D_3.
\end{align*}
The sum is
\[
N_{EDD}=D_1D_2E_3-D_1D_3E_2+D_1E_2D_3-D_3E_1D_2+E_1D_2D_3 .
\]
Right route, starting with $\sigma_1$ on $D_2D_1$:
\begin{align*}
E_3D_2D_1&\to E_3D_1D_2\to D_1E_3D_2-D_3E_1D_2+E_1D_3D_2\\
&\to D_1D_2E_3-D_1D_3E_2+D_1E_2D_3-D_3E_1D_2+E_1D_2D_3=N_{EDD}.
\end{align*}

\emph{$EED$.} Left route, starting with $\sigma_2$ on $E_3E_2$:
\begin{align*}
E_3E_2D_1&\to E_2E_3D_1\to E_2D_1E_3-E_2D_3E_1+E_2E_1D_3\\
&\to D_1E_2E_3-D_2E_1E_3+E_1D_2E_3-E_2D_3E_1+E_1E_2D_3=:N_{EED}.
\end{align*}
Right route, starting with $\sigma_3$ on $E_2D_1$:
\[
E_3E_2D_1\to E_3D_1E_2-E_3D_2E_1+E_3E_1D_2.
\]
We reduce each term:
\begin{align*}
E_3D_1E_2&\to D_1E_3E_2-D_3E_1E_2+E_1D_3E_2\to D_1E_2E_3-D_3E_1E_2+E_1D_3E_2,\\
-E_3D_2E_1&\to -D_2E_3E_1+D_3E_2E_1-E_2D_3E_1\to -D_2E_1E_3+D_3E_1E_2-E_2D_3E_1,\\
E_3E_1D_2&\to E_1E_3D_2\to E_1D_2E_3-E_1D_3E_2+E_1E_2D_3.
\end{align*}
The sum is again $N_{EED}$.

In all four cases the two routes have the same normal form, so all ambiguities are resolvable.
\end{proof}

\begin{theorem}\label{thm:pbw}
The normal words form a basis of $U(\g)$. In particular, for every content and bidegree, the dimension of the corresponding component of $U(\g)$ equals the number of normal words of that content and bidegree.
\end{theorem}

\begin{proof}
We apply Bergman's diamond lemma \cite[Theorem~1.2]{Ber}. Its hypotheses hold: $\kk$ is a commutative ring, $\mathcal S$ is a reduction system, and $\le$ is a semigroup partial order on $\langle X\rangle$, compatible with $\mathcal S$ and satisfying the descending chain condition. The theorem says that resolvability of all ambiguities (Proposition~\ref{prop:ambig}) is equivalent to the irreducible monomials forming a basis of $\kk\langle X\rangle/I$, where $I$ is the ideal generated by the differences $W_\sigma-f_\sigma$. Here $I=(\mathcal R)$. All rules are homogeneous in content and bidegree, which gives the second statement.
\end{proof}

A word of content a set $S$ is a sequence $X^{(1)}_{\sigma_1}X^{(2)}_{\sigma_2}\cdots X^{(n)}_{\sigma_n}$, where $(\sigma_1,\dots,\sigma_n)$ is an ordering of $S$ and $X^{(p)}\in\{D,E\}$. Such a word is normal if and only if at every descent $\sigma_p>\sigma_{p+1}$ the labels are $(X^{(p)},X^{(p+1)})=(D,E)$.

\section{Counting normal words}\label{sec:count}

For a word $w$ let $\omega(w)=y^{\#D}z^{\#E}$. For $n\ge 0$, let $u_n$ be the sum of $\omega(w)$ over the normal words of content $[n]$; normality depends only on the relative order of indices. By Theorem~\ref{thm:pbw},
\[
\Hs_{U(\g)}(x)=\sum_{n\ge0}u_n\frac{x^n}{n!}.
\]
Call a word $c$ of content a set a \emph{chain} if each of its adjacent pairs is bad. The empty word and words of length $1$ are chains.

\begin{lemma}\label{lem:chain}
Let $S$ be a set with $|S|=k$. The chains of content $S$ are exactly the words whose indices are in decreasing order and whose labels have the form $E\cdots ED\cdots D$. Their total weight is
\[
h_k=\sum_{a=0}^ky^az^{k-a}.
\]
\end{lemma}

\begin{proof}
A pair $X_jY_i$ is bad if and only if $j>i$ and $(X,Y)\neq(D,E)$. So in a chain the indices decrease, and no $D$ is immediately followed by an $E$. Conversely, every such word is a chain. There is one such word for each number $a\in\{0,\dots,k\}$ of $D$'s.
\end{proof}

\begin{proposition}\label{prop:involution}
For $n\ge1$,
\[
\sum_{k=0}^n(-1)^k\binom nk h_k\,u_{n-k}=0.
\]
\end{proposition}

\begin{proof}
Let $\Omega_n$ be the set of pairs $(c,w)$ where $c$ is a chain, $w$ is a normal word, and the contents of $c$ and $w$ partition $[n]$. Give $(c,w)$ sign $(-1)^{|c|}$ and weight $\omega(c)\omega(w)$. Grouping by the content of $c$ and using Lemma~\ref{lem:chain}, the signed weighted count of $\Omega_n$ is the left-hand side. It therefore suffices to construct a fixed-point-free involution $\Phi$ of $\Omega_n$ that preserves weight and reverses sign.

For $(c,w)\in\Omega_n$:
\begin{itemize}
\item[(I)] If $w$ is nonempty and either $c$ is empty or the pair $(\mathrm{last}(c),\mathrm{first}(w))$ is bad, let $\Phi(c,w)$ be obtained by moving the first letter of $w$ to the end of $c$.
\item[(II)] Otherwise, let $\Phi(c,w)$ be obtained by moving the last letter of $c$ to the front of $w$.
\end{itemize}

In case (II), $c$ is nonempty. Otherwise $w$ would be nonempty, since $n\ge1$, and case (I) would apply.

\emph{$\Phi$ maps $\Omega_n$ to itself.} In case (I), the new chain is a chain, because the only new adjacent pair is bad. The new word is a suffix of a normal word, hence normal. In case (II), the new chain is a prefix of a chain. The only new adjacent pair in the new word is $(\mathrm{last}(c),\mathrm{first}(w))$, which is not bad by the case distinction (or $w$ is empty), so the new word is normal.

\emph{$\Phi$ is an involution.} Suppose (I) moved $f=\mathrm{first}(w)$, giving $(cf,w')$. If $w'$ is nonempty, then $(f,\mathrm{first}(w'))$ is an adjacent pair of the normal word $w$, hence not bad. So (I) does not apply to $(cf,w')$, and (II) moves $f$ back.

Suppose (II) moved $\ell=\mathrm{last}(c)$, giving $(c',\ell w)$ with $c=c'\ell$. Then $\ell w$ is nonempty, and either $c'$ is empty or $(\mathrm{last}(c'),\ell)$ is an adjacent pair of the chain $c$, hence bad. So (I) applies and moves $\ell$ back.

Finally, $\Phi$ changes $|c|$ by $\pm1$ and preserves the multiset of letters. Hence it reverses the sign, preserves the weight, and has no fixed points.
\end{proof}

\begin{corollary}\label{cor:Useries}
Let $H(x)=\sum_{k\ge0}h_k\frac{x^k}{k!}$. Then $H(-x)\,\Hs_{U(\g)}(x)=1$, and
\[
(y-z)H(x)=y\,e^{xy}-z\,e^{xz}.
\]
\end{corollary}

\begin{proof}
The coefficient of $x^n/n!$ in $H(-x)\Hs_{U(\g)}(x)$ is $\sum_k(-1)^k\binom nkh_ku_{n-k}$. This is $1$ for $n=0$, since $h_0=u_0=1$, and $0$ for $n\ge1$ by Proposition~\ref{prop:involution}. The second identity follows from $(y-z)h_k=y^{k+1}-z^{k+1}$.
\end{proof}

\section{The Hilbert series of \texorpdfstring{$\g$}{g} and the proof of Theorem~\ref{thm:main}}\label{sec:hilbert}

\begin{proposition}\label{prop:gseries}
$\Hs_{\g}(x)=x(y+z)-\log\big(1-\mathcal Y(x)\big)$ and $\Hs_{\kr}(x)=xz-\log\big(1-\mathcal Y(x)\big)=\Hs_{\ktop}(x)$.
\end{proposition}

\begin{proof}
By (F4) and Corollary~\ref{cor:Useries}, $\Hs_{\g}=\log\Hs_{U(\g)}=-\log H(-x)$. We claim that
\begin{equation}\label{eq:key}
H(-x)=e^{-x(y+z)}\big(1-\mathcal Y(x)\big).
\end{equation}
First, the coefficient of $x^w/w!$ in $H(x)-e^{xy}-e^{xz}+1$ is
\begin{itemize}
\item $1-1-1+1=0$ for $w=0$;
\item $(y+z)-y-z=0$ for $w=1$;
\item $h_w-y^w-z^w=\sum_{a=1}^{w-1}y^az^{w-a}=\sum_{s=2}^wy^{s-1}z^{w-s+1}$ for $w\ge2$.
\end{itemize}
Hence $\mathcal Y=H(x)-e^{xy}-e^{xz}+1$. Using Corollary~\ref{cor:Useries},
\[
(y-z)(1-\mathcal Y)=(y-z)(e^{xy}+e^{xz})-(y\,e^{xy}-z\,e^{xz})=y\,e^{xz}-z\,e^{xy}.
\]
On the other hand,
\[
(y-z)\,e^{x(y+z)}H(-x)=e^{x(y+z)}\big(y\,e^{-xy}-z\,e^{-xz}\big)=y\,e^{xz}-z\,e^{xy}.
\]
Since $y-z$ is not a zero divisor in $\Q[y,z][[x]]$, \eqref{eq:key} follows. Therefore
\[
\Hs_\g=-\log H(-x)=x(y+z)-\log(1-\mathcal Y).
\]
By Proposition~\ref{prop:k}(3), $\g\cong\kr\oplus\af_D$ as graded spaces, and $\Hs_{\af_D}=xy$. Hence $\Hs_\kr=\Hs_\g-xy$. Comparing with Corollary~\ref{cor:ktopseries} gives the second statement.
\end{proof}

\begin{proof}[Proof of Theorem~\ref{thm:main}]
By Proposition~\ref{prop:gseries}, $\dim\kr_{[n]}^{(a,b)}=\dim(\ktop)_{[n]}^{(a,b)}$ for all $n,a,b$. Summing over $(a,b)$ gives $\dim\kr_{[n]}=\dim\ktop_{[n]}$. By Lemma~\ref{lem:A}, $(\ker\theta)_{[n]}=0$ for all $n$. By Proposition~\ref{prop:polar}, $\ker\theta=0$.
\end{proof}

\section{Corollaries}\label{sec:cor}

\begin{corollary}\label{cor:structure}
Let $W=C+\{a,b\}$ with $C=\{j_1,\dots,j_r\}$. The element
\[
Y_W=\ad(E_{j_1})\cdots\ad(E_{j_r})[D_a,E_b]\in\h
\]
depends only on $W$. The $Y_W$, with $|W|\ge2$, freely generate a free Lie subalgebra $\Lf\subset\h$. Moreover $\h=\Lf\rtimes\af_E$, where $\af_E=\operatorname{span}(E_i)$ is abelian and acts by $[E_j,Y_W]=Y_{W+j}$. The multilinear Hilbert series of $\h$ is $xz-\log(1-\mathcal Y_2(x))$.
\end{corollary}

\begin{proof}
By Theorem~\ref{thm:main}, $\theta\colon\lie\to\h$ is an isomorphism. By Lemma~\ref{lem:J0}, $\lie$ is the excess-$0$ part of $\ktop$. By Lemma~\ref{lem:B} and the proof of Corollary~\ref{cor:ktopseries}, that part is $\Lf(Y_{\cdot,2})\rtimes\af_E$, with $Y_{W,2}\mapsto\ad(E_{C})G_{\{a,b\}}$. Finally, $\theta$ maps $\ad(E_{C})F_{ab}$ to $Y_W$.
\end{proof}

\begin{corollary}\label{cor:free}
The elements $[D_i,E_j]$, $i\le j$, freely generate a free Lie subalgebra of $\g$.
\end{corollary}

\begin{proof}
These elements are $Y_{\{i,j\}}$, a subset of the free generators in Corollary~\ref{cor:structure}. A subset of free generators of a free Lie algebra freely generates a free Lie subalgebra, since this subalgebra is a retract.
\end{proof}

For $m\ge0$ let $A_m(t)=\sum_{\sigma\in\mathfrak S_m}t^{\operatorname{des}\sigma}=\sum_dA(m,d)t^d$, where $\operatorname{des}\sigma=\#\{p:\sigma_p>\sigma_{p+1}\}$ and $A_0=1$.

\begin{lemma}\label{lem:euler}
Let $\mathcal A(x,t)=\sum_{m\ge0}A_m(t)\frac{x^m}{m!}\in\Q[t][[x]]$. Then $\big(t-e^{(t-1)x}\big)\,\mathcal A(x,t)=t-1$.
\end{lemma}

\begin{proof}
We first show that, for every $m\ge0$,
\begin{equation}\label{eq:carlitz}
\sum_{q\ge0}(q+1)^mt^q=\frac{A_m(t)}{(1-t)^{m+1}}\qquad\text{in }\Q[[t]].
\end{equation}
The number $(q+1)^m$ counts maps $f\colon[m]\to\{0,\dots,q\}$. To such an $f$ associate the permutation $\sigma$ listing $[m]$ by increasing value of $f$, ties broken by increasing index. Then $f(\sigma_1)\le\cdots\le f(\sigma_m)$, with strict inequality at every descent of $\sigma$. Conversely, a permutation $\sigma$ and a weakly increasing sequence $f(\sigma_1)\le\cdots\le f(\sigma_m)$ that is strict at the descents of $\sigma$ determine $f$, and $f$ gives back $\sigma$.

For fixed $\sigma$ with $d$ descents, subtracting from $f(\sigma_p)$ the number of descents at positions $<p$ gives a bijection onto the weakly increasing sequences with values in $\{0,\dots,q-d\}$. There are $\binom{q-d+m}{m}$ of these, and
\[
\sum_q\binom{q-d+m}{m}t^q=\frac{t^d}{(1-t)^{m+1}}.
\]
Summing over $\sigma$ proves \eqref{eq:carlitz}.

Now work in $\Q[[t,x]]$. By \eqref{eq:carlitz},
\[
\mathcal A(x,t)=\sum_m\frac{x^m}{m!}(1-t)^{m+1}\sum_q(q+1)^mt^q=(1-t)\sum_qt^qe^{(q+1)(1-t)x}=\frac{(1-t)e^{(1-t)x}}{1-te^{(1-t)x}}.
\]
Here the inner sums are rearranged legitimately in $\Q[[t,x]]$, and the denominator has constant term $1$. Multiplying by $1-te^{(1-t)x}$ and then by $-e^{(t-1)x}$ gives $(t-e^{(t-1)x})\mathcal A=t-1$ in $\Q[[t,x]]$, hence in $\Q[t][[x]]$.
\end{proof}

\begin{corollary}\label{cor:eulerian}
For $n\ge2$ and $0\le a\le n$, $\dim\g_{[n]}^{(a,n-a)}=A(n-1,a-1)$, with the convention $A(n-1,-1)=A(n-1,n-1)=0$. In particular $\dim\g_{[n]}=(n-1)!$. The same holds for $\kr$ in place of $\g$.
\end{corollary}

\begin{proof}
Every letter carries exactly one index, so in content $[n]$ the bidegree is $(a,n-a)$. We may therefore set $z=1$ without losing information. Let $Q=(1-\mathcal Y)|_{z=1}\in\Q[y][[x]]$; it has constant term $1$. By the identity $(y-z)(1-\mathcal Y)=y\,e^{xz}-z\,e^{xy}$ from the proof of Proposition~\ref{prop:gseries},
\[
(y-1)Q=y\,e^{x}-e^{xy},\qquad (y-1)Q'=y\,e^{x}-y\,e^{xy},
\]
where $'=d/dx$. Put $\mathcal A=\mathcal A(x,y)$. Lemma~\ref{lem:euler} gives $(y-e^{(y-1)x})(\mathcal A-1)=e^{(y-1)x}-1$. Therefore
\[
(y-1)\,y(\mathcal A-1)Q=y\,e^x(\mathcal A-1)\big(y-e^{(y-1)x}\big)=y\,e^{xy}-y\,e^x=-(y-1)Q'.
\]
Since $y-1$ is not a zero divisor, $-Q'/Q=y(\mathcal A-1)$, that is,
\[
-\frac{d}{dx}\log Q=\sum_{m\ge1}yA_m(y)\frac{x^m}{m!}.
\]
Integrating, using $\log Q(0)=0$, and recalling $\Hs_\g=x(y+z)-\log(1-\mathcal Y)$, the coefficient of $x^n/n!$ in $\Hs_\g|_{z=1}$ is $yA_{n-1}(y)$ for $n\ge2$. This is the first claim, and $A_{n-1}(1)=(n-1)!$ gives the second. For $\kr$, note that $\Hs_\kr$ and $\Hs_\g$ differ only in degree $n=1$.
\end{proof}

\begin{remark}[Koszulness]\label{rem:koszul}
This remark is not used elsewhere in the paper. Fix $N\in\N$, and let $\g_N$ be the Lie algebra with the presentation \eqref{eq:g}, but with indices restricted to $[N]$. Killing the generators with index $>N$ shows that $\g_N$ is the subalgebra of $\g$ generated by the $D_i,E_i$ with $i\le N$, and that it contains all components of $\g$ whose content is supported in $[N]$.

With $V=\kk^N$ and $U=\kk d\oplus\kk e$ as in Section~\ref{sec:polar}, we have $U(\g_N)=T(V\otimes U)/(R)$, where $R\subset(V\otimes U)^{\otimes2}\cong V^{\otimes2}\otimes U^{\otimes2}$ is the space $\Lambda^2V\otimes\Sym^2U$ of tensors antisymmetric in $V$ and symmetric in $U$. Indeed, $[D_i,D_j]$, $[E_i,E_j]$ and $[D_i,E_j]-[D_j,E_i]$ correspond to $(v_iv_j-v_jv_i)\otimes dd$, $(v_iv_j-v_jv_i)\otimes ee$ and $(v_iv_j-v_jv_i)\otimes(de+ed)$, and both spaces have dimension $3\binom N2$.

Since $\operatorname{char}\kk\ne2$, $V^{\otimes2}=\Sym^2V\oplus\Lambda^2V$ (and likewise for $U$), hence $R^\perp=\Sym^2V^*\otimes U^{*\otimes2}+V^{*\otimes2}\otimes\Lambda^2U^*$. This subspace is exactly the degree-$2$ kernel of the natural surjection onto the Segre product $\Lambda(V^*)\circ\Sym(U^*)$ \cite[1.2]{BF}: in degree $2$ that map is the tensor product of the quotient maps onto $\Lambda^2(V^*)$ and $\Sym^2(U^*)$. The Segre product is generated in degree $1$ and is $2$-related by \cite[Prop.~3(b)]{BF}. Here the Segre product uses ordinary tensor-product multiplication, not a super tensor product. All these algebras have finitely many degree-$1$ generators and finitely generated quadratic relation ideals, as required by the conventions of \cite[\S1]{BF}. Therefore the dual algebra $U(\g_N)^\circ$ in the sense of \cite[1.14]{BF} is isomorphic to $\Lambda(V^*)\circ\Sym(U^*)$.

The polynomial algebra $\Sym(V)$ is Koszul \cite[1.17(b)]{BF}, hence so is its dual $\Lambda(V^*)$ \cite[1.16(10)]{BF}; likewise $\Sym(U^*)$ is Koszul. By \cite[Theorem~4(b)]{BF} the Segre product is Koszul, and by \cite[1.16(10)]{BF} so is $U(\g_N)$.

Give $v_i^*$ content $i$ and $d^*,e^*$ the bidegrees of $d,e$. For $0\le n\le N$, the content-$[n]$ part of $\Lambda^n(V^*)$ is one-dimensional, so the content-$[n]$ part of $\Lambda(V^*)\circ\Sym(U^*)$ is $\Sym^n(U^*)$, of weight $h_n$. Taking $N\ge n$ in each degree, the identity $H(-x)\Hs_{U(\g)}(x)=1$ of Corollary~\ref{cor:Useries} has the shape of the numerical relation between a Koszul algebra and its dual \cite[1.16(5)]{BF}; we have proved it directly. Finally, Dotsenko, Shadrin and Tamaroff prove that $\mathrm{tHyperCom}$ is Koszul \cite[Prop.~2.4.3]{DST}. In view of Section~\ref{sec:prior}, that statement concerns the multilinear species associated with the enveloping algebra of the ideal $\kr$ in its generators $G_T$ (with weight given by the number of letters), not directly the ordinary algebra $U(\g)$ in the generators $D_i,E_i$. We do not infer Koszulness of an ordinary algebra from this species statement; the finite-generator argument above establishes the assertion about $U(\g_N)$.
\end{remark}

\section{Relation to prior work}\label{sec:prior}

In this section we explain how the dimension results of Sections~\ref{sec:pbw}--\ref{sec:cor} follow from the work of Losev--Manin \cite{LM} and Dotsenko--Shadrin--Tamaroff \cite{DST}. We also distinguish results not explicitly stated there. Numbered statements of \cite{LM} and \cite{DST} refer to arXiv version~2 of \cite{LM} and arXiv version~3 of \cite{DST}. We use homology with characteristic-$0$ coefficients; the algebraic conventions of \cite[\S2.1]{DST} are in characteristic~$0$.

\subsection*{Losev--Manin spaces}
For $n\ge1$, Losev and Manin introduced a smooth proper variety $L_n$ of dimension $n-1$ \cite[Thm.~2.2]{LM}, which they identify with the toric variety associated with the permutohedron \cite[\S2]{LM}. In \cite[Def.~2.4.1]{DST} it is denoted $\mathrm{LM}(n)$ and described as the toric variety of the fan of Weyl chambers of type $A$. By \cite[Thm.~2.3]{LM}, $L_n$ has no odd cohomology, and
\[
1+\sum_{n\ge1}p_n(q)\frac{y^n}{n!}=\frac{q-1}{q-e^{(q-1)y}},\qquad p_n(q)=\sum_i\dim H^{2i}(L_n)\,q^i.
\]
Both this series and $\mathcal A(y,q)$ of Lemma~\ref{lem:euler} satisfy $(q-e^{(q-1)y})F=q-1$, and $\Q[q][[y]]$ is an integral domain, so $p_n(q)=A_n(q)$: the Betti numbers of $L_n$ are Eulerian numbers. The strata of $L_n$ are indexed by ordered set partitions of $[n]$, and inclusions of boundary strata make the homology a twisted associative algebra, i.e.\ an associative monoid in species \cite[\S2.4]{DST}. Dotsenko, Shadrin and Tamaroff showed \cite[Prop.~2.4.2]{DST} that this algebra is the twisted associative algebra $\mathrm{tHyperCom}$. It is generated by fully symmetric elements $m_t$ of arity $t\ge1$ and homological degree $2t-2$, subject to
\[
\sum_{\substack{A=I\sqcup J\\ i\in I,\ j\in J}}\big(m_Im_J-m_Jm_I\big)=0\qquad(i\ne j\in A).
\]
Their proof combines the relations established in \cite{LM}, the fact that the total Betti number of $L_n$ is the number $n!$ of chambers, and a Gr\"obner basis argument giving $\dim\mathrm{tHyperCom}(n)\le n!$. They credit the monomial ordering in that argument to an unpublished internship report of P.~Laubie, which we have not seen. All generators have even degree and trivial symmetric-group action, so no Koszul signs occur.

\subsection*{The dictionary}
We keep the singleton generators: $G_{\{i\}}=E_i$ corresponds to $m_1$, the fundamental class of $L_1$, which is a point. There is no generator $G_\emptyset$ or $m_0$. In arity $0$ we adjoin only the algebra unit, so $U(\kr)_\emptyset=\kk$; the constant $1$ in the generating series is this formal unit and makes no assertion about a geometric space $L_0$. In every ordered set partition below the blocks are nonempty.

By (F3), $U(\kr)=\kk\langle G_T\rangle/(r_{i,j,S})$. Fix $n\ge1$. The component of $U(\kr)$ of content $[n]$ is the span of the words $G_{T_1}\cdots G_{T_p}$, where $(T_1,\dots,T_p)$ runs over the ordered set partitions of $[n]$. One quotients by the span of the elements $u\,r_{i,j,S}\,v$ in which $u$, $v$ are words and all contents are disjoint sets with union $[n]$. This is, by definition, arity $n$ of the twisted associative algebra generated by the symbols $G_T$ (with trivial symmetric-group action) subject to the relators $r_{i,j,S}$ for sets $S\sqcup\{i,j\}$. For a set $S$, \eqref{eq:r} reads
\[
r_{i,j,S}=\sum_{\substack{A=I\sqcup J\\ i\in I,\ j\in J}}\big(G_IG_J-G_JG_I\big),\qquad A=S\sqcup\{i,j\},
\]
in $U(\kr)$, which is the relator of $\mathrm{tHyperCom}$ under $G_T\leftrightarrow m_T$. A word with $p$ letters has $D$-degree $n-p$, while $m_{T_1}\cdots m_{T_p}$ has homological degree $\sum_k(2|T_k|-2)=2(n-p)$. Therefore, by \cite[Prop.~2.4.2]{DST} and \cite[Thm.~2.3]{LM},
\begin{equation}\label{eq:LMdims}
\dim U(\kr)_{[n]}^{(a,n-a)}=\dim H_{2a}(L_n;\kk)=\dim H^{2a}(L_n)=A(n,a),\qquad \dim U(\kr)_{[n]}=n!.
\end{equation}
(Homology and cohomology with field coefficients have the same dimensions.) For $n=2$, the relator is $[G_i,G_j]=0$, matching $m_{\{i\}}m_{\{j\}}=m_{\{j\}}m_{\{i\}}$. The identification is one of species, compatible with products on disjoint index sets; we are not identifying the ordinary algebra $U(\kr)$, including its repeated-index components, with a direct sum of unlabelled homology spaces. It involves only a doubling of the grading and no signs. We checked it directly for $n\le6$ by computing $U(\kr)_{[n]}$ from our relators over $\Q$ (script \texttt{prior\_work\_check.py}).

\subsection*{Consequences}
By Proposition~\ref{prop:k}(3) and (F4), $\Hs_{U(\g)}=\exp(\Hs_\kr+xy)=e^{xy}\,\Hs_{U(\kr)}$. In content $[n]$ the bidegree is $(a,n-a)$, so we may set $z=1$, and \eqref{eq:LMdims} gives
\[
\Hs_{U(\g)}\big|_{z=1}=e^{xy}\,\frac{y-1}{y-e^{(y-1)x}}=\frac{y-1}{y\,e^{-xy}-e^{-x}}.
\]
This is Corollary~\ref{cor:Useries} at $z=1$; we confirmed the identity with \texttt{sympy}. Taking logarithms gives Proposition~\ref{prop:gseries} and Corollary~\ref{cor:eulerian}. Thus the dimension count used in the proof of Theorem~\ref{thm:main}, namely $\dim\kr_{[n]}^{(a,b)}=\dim(\ktop)_{[n]}^{(a,b)}$, could also be obtained by combining \cite{LM,DST} with Proposition~\ref{prop:k} and Corollary~\ref{cor:ktopseries}. Corollary~\ref{cor:eulerian} and the counts of Sections~\ref{sec:pbw}--\ref{sec:hilbert} should therefore be regarded as known, in an equivalent form.

Our derivation is independent. We use a quadratic Gr\"obner basis of $U(\g)$ in the letters $D_i,E_i$ (Section~\ref{sec:pbw}), whereas \cite{DST} use the generators $G_T$ of $U(\kr)$. We then count normal words by a sign-reversing involution (Section~\ref{sec:count}), whereas \cite{LM} count points of $L_n$ over finite fields.

\subsection*{Results not explicitly stated in \texorpdfstring{\cite{LM,DST}}{[LM, DST]}}
The works \cite{LM,DST} concern the homology algebras of the Losev--Manin spaces and their representation-theoretic applications. They do not explicitly formulate $\g$, $\kr$ and $\h$ with the presentations used here, or ask for the relations of $\h$. Lemma~\ref{lem:A}, Lemma~\ref{lem:B}, Theorem~\ref{thm:main} and Corollaries~\ref{cor:structure} and~\ref{cor:free} are not explicitly stated in \cite{LM,DST}; we were unable to find these statements elsewhere in the literature. This is not a claim that they cannot be deduced from earlier work.

\subsection*{A heuristic explanation of the link}
Since $\kk$ is infinite, a representation $\rho$ of $\g$ is the same as a family of operators $X_i=\rho(D_i)$, $Y_i=\rho(E_i)$ with $[X_i+\lambda Y_i,X_j+\lambda Y_j]=0$ for all $\lambda\in\kk$: expand in $\lambda$. For instance, let $A=\sum_iA_i\,dt_i$ be a matrix-valued $1$-form in formal variables $t_1,\dots,t_N$. Put $X_i=\partial/\partial t_i$ and let $Y_i$ be multiplication by $A_i$, acting on vector-valued formal power series. Then these relations say $\partial_iA_j=\partial_jA_i$ and $[A_i,A_j]=0$, i.e.\ that $d+\lambda A$ is a pencil of flat connections in the sense of \cite[\S3.1--3.2]{LM}. In this representation $\rho(\mathbf G_T)$ is multiplication by $\partial_{T-t}A_t$. If $A=dB$ \cite[Prop.~3.2]{LM}, this is $\partial_TB$. For $B$ as in \cite[(3.19)]{LM}, the value of $\partial_TB$ at $t=0$ is the ``top correlator'' indexed by $T$. We include this only as motivation; nothing above depends on it.

\section{Computational verification}\label{sec:computations}

The proof above is complete without computer assistance: Proposition~\ref{prop:ambig} is checked by hand. We nevertheless verified its ingredients and its conclusion independently by computer. All scripts are in Python~3 and are contained in the accompanying archive \texttt{code.zip}, together with their outputs.

Throughout this section, $k$ denotes the $E$-degree of a multilinear component of content $[n]$. For elements of $\kr$ or $\lie$ in the $G_T$-presentation, it equals the number of letters $G_T$. Modular computations use the prime $p=2^{31}-1$.

\begin{table}[ht]
\centering
\small
\begin{tabular}{@{}>{\raggedright\arraybackslash}p{0.34\linewidth}>{\raggedright\arraybackslash}p{0.2\linewidth}>{\raggedright\arraybackslash}p{0.1\linewidth}>{\raggedright\arraybackslash}p{0.27\linewidth}@{}}
\toprule
Check & Script & Range & Result\\
\midrule
$\dim\kr$, $\dim\lie$, $\dim\h$, $\dim\ker\theta$ in each component $([n],k)$ & \texttt{kmodel\_exact.py} (exact over $\Q$) & $n\le8$ & $\ker\theta=0$; $\dim\kr$ Eulerian\\
same, modular & \texttt{kmodel.py} & $n\le8$ & identical\\
$\g$ directly in the letters $D_i,E_i$ (independent model) & \texttt{gmodel.py} (mod $p$) & $n\le6$ & $\dim\g$ Eulerian; $\dim\h$ agrees\\
$\dim\ktop$ & \texttt{ktop.py} (mod $p$) & $n\le8$ & matches Cor.~\ref{cor:ktopseries}\\
degree-$3$ Gr\"obner check: $\dim(X\mathcal R+\mathcal RX)_3=\#$ non-normal words & \texttt{pbw\_check.py} (exact) & indices $1..3$ & $104=104$\\
explicit ambiguity resolution (Prop.~\ref{prop:ambig}) & \texttt{ambiguities.py} (exact) & all $4$ & resolvable\\
normal-word counts, then $\log$ & \texttt{pbw\_count.py} (exact) & $n\le10$ & equals Prop.~\ref{prop:gseries}\\
identity \eqref{eq:key} and Eulerian numbers & \texttt{egf\_check.py} (sympy) & $n\le10$ & identical\\
injectivity via projection (below) & \texttt{kproj.py} (mod $p$) & $n=9$ & $\ker\theta=0$\\
dictionary of Section~\ref{sec:prior}: relators, $\dim U(\kr)_{[n]}$ by degree, EGF & \texttt{prior\_work\_}\allowbreak\texttt{check.py} (exact, sympy) & $n\le6$ & agrees with \cite{LM,DST}\\
\bottomrule
\end{tabular}
\caption{Computer checks. The formulas for $\dim\kr$ and $\dim\lie$ are tabulated by \texttt{lformula.py}.}
\label{tab:checks}
\end{table}

\subsection*{Degree-$3$ Gr\"obner check}
For a quadratic algebra $\kk\langle X\rangle/(\mathcal R)$, the normal words of degree $3$ are linearly independent modulo $(\mathcal R)_3=X\mathcal R+\mathcal RX$ if and only if $\dim(X\mathcal R+\mathcal RX)=\#\{\text{words of length }3\text{ containing a bad subword}\}$. Indeed, the normal words always span modulo $(\mathcal R)_3$, by reduction.

This independence implies that every overlap ambiguity resolves. Indeed, the two reductions of an overlap differ by an element of $(\mathcal R)_3$, and their normal forms are then equal. Since an overlap involves at most three distinct indices and the rules depend only on the order of the indices, indices $\{1,2,3\}$ suffice.

The script \texttt{pbw\_check.py} computes this in exact rational arithmetic and finds $104=104$. As a sanity check, it reports the same equality for all eight natural letter orders. The companion script \texttt{sanity\_random.py} applies the same test to three random sets of nine quadratic relations in the same letters. The test fails in each case (for example $108\ne92$), which shows that it is not vacuous. Note that $\dim(X\mathcal R+\mathcal RX)\ge\#\{\text{non-normal words}\}$ always holds, since the normal words span.

\subsection*{Dimensions}
The script \texttt{pbw\_count.py} counts multilinear normal words by dynamic programming. For $n=1,\dots,10$ the totals are $2,5,16,65,326,1957,13700,109601,986410,9864101$. It then checks that the logarithm of their exponential generating function equals $x(y+z)-\log(1-\mathcal Y)$ through order $10$. The resulting $\dim\g_{[n]}^{(a,n-a)}$ are Eulerian numbers. For instance, for $n=10$ they are $1,502,14608,88234,156190,88234,14608,502,1$.

\subsection*{Injectivity in degree $9$}
Let $F_{>0}=\bigoplus_{q>0}F_q$ and let $\pi\colon F\to F_{>0}$ be the projection along $F_0$. Since $\h\cong(F_0+I)/I$ and $F=F_0\oplus F_{>0}$,
\[
\dim(\kr/\h)_{[n],k}=\dim(F_{>0})_{[n],k}-\operatorname{rank}\pi(I_{[n],k}).
\]
The script \texttt{kproj.py} computes $\pi(I)$ recursively modulo $p$. Free Lie elements are encoded by their coefficients on associative words starting with the letter containing the minimal index; these coordinates are valid over every field. For $n=9$ it finds, for $k=1,\dots,9$:
\[
\dim(\kr/\h)_{[9],k}=1,\ 247,\ 4293,\ 15619,\ 8059,\ 457,\ 1,\ 0,\ 0.
\]
These are exactly $\dim\kr_{[9],k}-\dim\lie_{[9],k}$, computed from Corollary~\ref{cor:ktopseries} and Proposition~\ref{prop:gseries}:
\begin{itemize}
\item $\dim\kr_{[9],k}=1,247,4293,15619,15619,4293,247,1,0$;
\item $\dim\lie_{[9],k}=0,0,0,0,7560,3836,246,1,0$.
\end{itemize}
The run took $554$\,s with a peak resident memory of $403$\,MB.

This modular computation is rigorous over $\Q$. The space $I_{[n],k}$ is spanned by integral vectors: the ad-words of letters applied to the integral relators $r_{i,j,S}$. The modular computation determines the span of their reductions modulo $p$, that is, $\pi(I)$ over $\mathbb F_p$. The rank over $\Q$ of a set of integral vectors is at least its rank modulo $p$. Hence
\[
\dim(\kr/\h)_\Q\le\dim(\kr/\h)_{\mathbb F_p}=\dim\kr-\dim\lie,
\]
and so $\dim\h_\Q\ge\dim\lie$. Since $\h$ is a quotient of $\lie$, equality holds, and $\theta$ is injective in content $[9]$. This argument uses only the dimension formulas, not Lemma~\ref{lem:A}.

A direct computation for $n=10$ was not attempted, because of its estimated time and memory requirements.

\subsection*{Reproducibility}
The archive contains all scripts, their outputs (\texttt{out\_*.txt}), and a README describing the command lines. The exact computations use Python's \texttt{fractions} module, and the series manipulations use \texttt{sympy}.

\subsection*{Acknowledgements}
The author used AI assistance in the literature search, in designing and running the computational experiments, in writing the accompanying verification code, and in drafting the manuscript. The human author is solely responsible for verifying all results, for the correctness of the proofs and computations, and for the final text.

\noindent\textit{Tools used:} Codex (OpenAI) and Claude (Anthropic), used as research and writing assistants. Specific version numbers were not recorded. AI is not listed as an author.

\providecommand{\bysame}{\leavevmode\hbox to3em{\hrulefill}\thinspace}
\begin{thebibliography}{10}

\bibitem{BF}
J\"orgen Backelin and Ralf Fr\"oberg, \emph{Koszul algebras, {V}eronese subrings and rings with linear resolutions}, Rev. Roumaine Math. Pures Appl. \textbf{30} (1985), no.~2, 85--97.

\bibitem{Ber}
George~M. Bergman, \emph{The diamond lemma for ring theory}, Adv. Math. \textbf{29} (1978), no.~2, 178--218.

\bibitem{DST}
Vladimir Dotsenko, Sergey Shadrin, and Pedro Tamaroff, \emph{Generalized cohomological field theories in the higher order formalism}, Comm. Math. Phys. \textbf{399} (2023), no.~3, 1439--1500, \href{https://doi.org/10.1007/s00220-022-04577-6}{doi:10.1007/s00220-022-04577-6}, \href{https://arxiv.org/abs/2112.06015}{arXiv:2112.06015}.

\bibitem{LM}
A. Losev and Yu. Manin, \emph{New moduli spaces of pointed curves and pencils of flat connections}, Michigan Math. J. \textbf{48} (2000), no.~1, 443--472, \href{https://doi.org/10.1307/mmj/1030132728}{doi:10.1307/mmj/1030132728}, \href{https://arxiv.org/abs/math/0001003}{arXiv:math/0001003}.

\bibitem{MO}
Terry Tao, \emph{Relations in a certain {L}ie algebra}, MathOverflow, question 464808, February 24, 2024, \url{https://mathoverflow.net/questions/464808}.

\bibitem{YCor}
YCor, comments on \cite{MO}, February 24--27, 2024 (UTC), \url{https://mathoverflow.net/posts/comments/1207321}. The reformulation and the question about the elements $[D_i,E_j]$ are in this first comment.

\end{thebibliography}

\end{document}
